\documentclass[a4paper,12pt]{book}

% ============================================================================
% PACKAGES ESSENTIELS
% ============================================================================
\usepackage[utf8]{inputenc}
\usepackage[T1]{fontenc}
\usepackage[french]{babel}
\usepackage[margin=2.5cm,top=3cm,bottom=3cm]{geometry}
\usepackage{xcolor}
\usepackage{titlesec}
\usepackage{fancyhdr}
\usepackage{hyperref}
\usepackage{listings}
\usepackage{mdframed}
\usepackage{tcolorbox}
\usepackage{longtable}
\usepackage{array}
\usepackage{graphicx}
\usepackage{tikz}
\usepackage{enumitem}
\usepackage{amsmath}
\usepackage{amssymb}
\usepackage{booktabs}
\usepackage{multirow}
\usepackage{colortbl}
\usepackage{fontawesome5}
\usepackage{csquotes}

% ============================================================================
% CONFIGURATION DES COULEURS
% ============================================================================
\definecolor{deepblue}{RGB}{0,51,102}
\definecolor{darkblue}{RGB}{0,102,204}
\definecolor{lightblue}{RGB}{102,178,255}
\definecolor{forestgreen}{RGB}{34,139,34}
\definecolor{darkgreen}{RGB}{0,100,0}
\definecolor{warningorange}{RGB}{255,140,0}
\definecolor{dangerred}{RGB}{220,20,60}
\definecolor{codegray}{RGB}{245,245,245}
\definecolor{codeborder}{RGB}{200,200,200}
\definecolor{architectureblue}{RGB}{25,25,112}
\definecolor{bestpracticegreen}{RGB}{0,128,0}
\definecolor{errorred}{RGB}{178,34,34}
\definecolor{definitionpurple}{RGB}{128,0,128}
\definecolor{analogybrown}{RGB}{139,69,19}

% ============================================================================
% CONFIGURATION DES LISTINGS (CODE)
% ============================================================================
\lstdefinestyle{javascript}{
    language=JavaScript,
    backgroundcolor=\color{codegray},
    commentstyle=\color{darkgreen}\itshape,
    keywordstyle=\color{deepblue}\bfseries,
    numberstyle=\tiny\color{gray},
    stringstyle=\color{dangerred},
    basicstyle=\ttfamily\footnotesize,
    breakatwhitespace=false,
    breaklines=true,
    captionpos=b,
    keepspaces=true,
    numbers=left,
    numbersep=5pt,
    showspaces=false,
    showstringspaces=false,
    showtabs=false,
    tabsize=2,
    frame=single,
    rulecolor=\color{codeborder},
    morekeywords={async,await,const,let,export,import,from,require,module}
}

\lstdefinestyle{sql}{
    language=SQL,
    backgroundcolor=\color{codegray},
    commentstyle=\color{darkgreen}\itshape,
    keywordstyle=\color{deepblue}\bfseries,
    numberstyle=\tiny\color{gray},
    stringstyle=\color{dangerred},
    basicstyle=\ttfamily\footnotesize,
    breaklines=true,
    numbers=left,
    numbersep=5pt,
    frame=single,
    rulecolor=\color{codeborder},
    morekeywords={CREATE,TABLE,INSERT,INTO,VALUES,SELECT,FROM,WHERE,AND,OR,NOT,NULL,DEFAULT,CHECK,UNIQUE,PRIMARY,KEY,FOREIGN,REFERENCES,ON,DELETE,CASCADE,UPDATE,RESTRICT,SET,EXTENSION,IF,NOT,EXISTS,UUID,GENERATE,ALWAYS,AS,TEXT,TIMESTAMP,WITH,TIME,ZONE,DEFAULT,CURRENT,DATE,RETURNING,ORDER,BY,LIMIT,OFFSET,JOIN,INNER,LEFT,RIGHT,FULL,OUTER,ON,GROUP,HAVING,COUNT,SUM,AVG,MAX,MIN,DISTINCT,ASC,DESC,LIKE,ILIKE,IN,BETWEEN,IS,TRUE,FALSE,AND,OR,NOT,UNION,INTERSECT,EXCEPT,ALL,ANY,SOME,CASE,WHEN,THEN,ELSE,END,CAST,EXTRACT,COALESCE,NULLIF,ARRAY,ROW,TYPE,ALTER,DROP,RENAME,ADD,COLUMN,CONSTRAINT,INDEX,CREATE,DROP,RENAME,ALTER,TABLE,INDEX,SCHEMA,DATABASE,USER,ROLE,GRANT,REVOKE,COMMIT,ROLLBACK,TRANSACTION,BEGIN,SAVEPOINT,RELEASE,LOCK,UNLOCK,SHARE,MODE,EXCLUSIVE,ACCESS,SHARE,UPDATE,NOWAIT,OF,ON,ALL,ANY,SOME,CURRENT,DATE,TIME,ZONE,INTERVAL,YEAR,MONTH,DAY,HOUR,MINUTE,SECOND,LOCAL,GLOBAL,SESSION,LEVEL,ISOLATION,READ,WRITE,ONLY,REPEATABLE,SERIALIZABLE,UNCOMMITTED,DEFERRED,IMMEDIATE,CONSTRAINTS,TRIGGER,FUNCTION,PROCEDURE,LANGUAGE,RETURNS,LANGUAGE,SQL,PLPGSQL,DECLARE,BEGIN,EXCEPTION,RETURN,PERFORM,RAISE,NOTICE,WARNING,ERROR,ASSERT,LOG,INFO,DEBUG}
}

\lstdefinestyle{bash}{
    language=bash,
    backgroundcolor=\color{codegray},
    commentstyle=\color{darkgreen}\itshape,
    keywordstyle=\color{deepblue}\bfseries,
    numberstyle=\tiny\color{gray},
    stringstyle=\color{dangerred},
    basicstyle=\ttfamily\footnotesize,
    breaklines=true,
    numbers=left,
    numbersep=5pt,
    frame=single,
    rulecolor=\color{codeborder},
    morekeywords={npm,node,cd,ls,mkdir,rm,cp,mv,git,git,init,add,commit,push,pull,clone,branch,checkout,merge,rebase,stash,pop,log,status,diff,show,remote,fetch,tag,config,help,version,install,update,uninstall,list,search,info,view,edit,remove,clean,cache,doctor,audit,fix,run,test,start,stop,restart,build,deploy,serve,dev,prod,env,set,unset,export,echo,cat,less,more,head,tail,grep,find,locate,which,where,type,man,help,history,alias,unalias,source,exec,exit,logout,sudo,su,chown,chgrp,chmod,chroot,mount,umount,df,du,free,top,ps,kill,pkill,killall,nice,renice,nohup,screen,tmux,ssh,scp,rsync,wget,curl,tar,gzip,gunzip,bzip2,bunzip,xz,unzip,zip,unrar,7z,sha256sum,md5sum,base64,openssl,python,python3,pip,pip3,javac,java,gcc,g++,make,cmake,npm,yarn,pnpm,bun}
}

\lstdefinestyle{html}{
    language=HTML,
    backgroundcolor=\color{codegray},
    commentstyle=\color{darkgreen}\itshape,
    keywordstyle=\color{deepblue}\bfseries,
    numberstyle=\tiny\color{gray},
    stringstyle=\color{dangerred},
    basicstyle=\ttfamily\footnotesize,
    breaklines=true,
    numbers=left,
    numbersep=5pt,
    frame=single,
    rulecolor=\color{codeborder}
}

\lstdefinestyle{css}{
    language=CSS,
    backgroundcolor=\color{codegray},
    commentstyle=\color{darkgreen}\itshape,
    keywordstyle=\color{deepblue}\bfseries,
    numberstyle=\tiny\color{gray},
    stringstyle=\color{dangerred},
    basicstyle=\ttfamily\footnotesize,
    breaklines=true,
    numbers=left,
    numbersep=5pt,
    frame=single,
    rulecolor=\color{codeborder}
}

% ============================================================================
% ENVIRONNEMENTS PERSONNALISÉS
% ============================================================================
\newtcolorbox{errorbox}[1][]{
    colback=red!5!white,
    colframe=dangerred,
    title=\textbf{\faExclamationTriangle\ Erreur},
    fonttitle=\bfseries,
    #1
}

\newtcolorbox{warningbox}[1][]{
    colback=orange!10!white,
    colframe=warningorange,
    title=\textbf{\faExclamation\ Attention},
    fonttitle=\bfseries,
    #1
}

\newtcolorbox{bestpracticebox}[1][]{
    colback=green!5!white,
    colframe=bestpracticegreen,
    title=\textbf{\faCheckCircle\ Bonne Pratique},
    fonttitle=\bfseries,
    #1
}

\newtcolorbox{definitionbox}[1][]{
    colback=purple!5!white,
    colframe=definitionpurple,
    title=\textbf{\faBook\ Définition},
    fonttitle=\bfseries,
    #1
}

\newtcolorbox{architecturebox}[1][]{
    colback=blue!5!white,
    colframe=architectureblue,
    title=\textbf{\faCogs\ Architecture},
    fonttitle=\bfseries,
    #1
}

\newtcolorbox{analogybox}[1][]{
    colback=brown!5!white,
    colframe=analogybrown,
    title=\textbf{\faLightbulb\ Analogie},
    fonttitle=\bfseries,
    #1
}

\newtcolorbox{securitybox}[1][]{
    colback=red!10!white,
    colframe=darkred,
    title=\textbf{\faShieldAlt\ Sécurité},
    fonttitle=\bfseries,
    #1
}

% ============================================================================
% CONFIGURATION DES TITRES
% ============================================================================
\titleformat{\chapter}[display]
{\normalfont\huge\bfseries\color{deepblue}}
{\chaptertitlename\ \thechapter}{20pt}{\Huge}

\titleformat{\section}
{\normalfont\Large\bfseries\color{darkblue}}
{\thesection}{1em}{}

\titleformat{\subsection}
{\normalfont\large\bfseries\color{lightblue}}
{\thesubsection}{1em}{}

% ============================================================================
% CONFIGURATION DES EN-TÊTES ET PIEDS DE PAGE
% ============================================================================
\pagestyle{fancy}
\fancyhf{}
\fancyhead[L]{\leftmark}
\fancyhead[R]{\thepage}
\fancyfoot[C]{\textit{Rapport Pédagogique - Blog Multi-Rôle}}
\renewcommand{\headrulewidth}{0.5pt}
\renewcommand{\footrulewidth}{0.5pt}

% ============================================================================
% CONFIGURATION DES HYPERLIENS
% ============================================================================
\hypersetup{
    colorlinks=true,
    linkcolor=deepblue,
    filecolor=magenta,
    urlcolor=darkblue,
    citecolor=darkgreen,
    pdfauthor={Université},
    pdftitle={Rapport Pédagogique - Blog Multi-Rôle},
    pdfsubject={Développement Backend Node.js + PostgreSQL}
}

% ============================================================================
% MÉTADONNÉES DU DOCUMENT
% ============================================================================
\title{\textbf{\Huge Rapport Pédagogique Complet}}
\subtitle{\Large Développement d'un Blog Multi-Rôle avec Node.js, Express et PostgreSQL}
\author{Licence 1 Informatique\\Département Génie Informatique}
\date{\today}

% ============================================================================
% DÉBUT DU DOCUMENT
% ============================================================================
\begin{document}

% ============================================================================
% PAGE DE GARDE
% ============================================================================
\begin{titlepage}
    \centering
    \vspace*{2cm}
    {\Huge\bfseries\color{deepblue} Rapport Pédagogique Complet\par}
    \vspace{1cm}
    {\Large\bfseries\color{darkblue} Développement d'un Blog Multi-Rôle\par}
    \vspace{0.5cm}
    {\large\itshape avec Node.js, Express et PostgreSQL\par}
    \vspace{2cm}
    
    \begin{tcolorbox}[colback=white,colframe=deepblue,width=0.8\textwidth,center,title=Stack Technologique]
        \begin{itemize}
            \item \textbf{Backend} : Node.js + Express
            \item \textbf{Base de données} : PostgreSQL
            \item \textbf{Driver} : pg (sans ORM)
            \item \textbf{Frontend} : HTML + CSS + JavaScript Vanilla
            \item \textbf{Authentification} : JWT
            \item \textbf{Sécurité} : bcrypt
        \end{itemize}
    \end{tcolorbox}
    
    \vspace{2cm}
    {\large Licence 1 Informatique\par}
    {\large Département Génie Informatique\par}
    \vspace{1cm}
    {\large \today\par}
\end{titlepage}

% ============================================================================
% TABLE DES MATIÈRES
% ============================================================================
\tableofcontents
\newpage

% ============================================================================
% CHAPITRE 1 : INTRODUCTION GÉNÉRALE
% ============================================================================
\chapter{Introduction Générale}

\section{Qu'est-ce qu'une Application Web ?}

Une application web est un logiciel qui fonctionne dans un navigateur web et communique avec un serveur via le protocole HTTP. Contrairement aux applications desktop traditionnelles qui sont installées localement sur un ordinateur, les applications web sont accessibles depuis n'importe quel appareil connecté à Internet.

\begin{definitionbox}
    \textbf{Application Web} : Programme informatique accessible via un navigateur web, utilisant les technologies web standard (HTML, CSS, JavaScript) et communiquant avec un serveur distant pour stocker et traiter les données.
\end{definitionbox}

\subsection{Architecture Client-Serveur}

L'architecture web moderne repose sur la séparation entre le \textbf{client} (frontend) et le \textbf{serveur} (backend) :

\begin{itemize}
    \item \textbf{Frontend (Client)} : L'interface utilisateur que l'on voit dans le navigateur. Il est responsable de l'affichage, de l'interaction avec l'utilisateur et de l'envoi des requêtes au serveur.
    \item \textbf{Backend (Serveur)} : Le "cerveau" de l'application qui traite les requêtes, exécute la logique métier et communique avec la base de données.
    \item \textbf{Base de données} : Le système de stockage persistant qui conserve les données de l'application (utilisateurs, articles, commentaires, etc.).
\end{itemize}

\begin{analogybox}
    \textbf{Analogie du Restaurant} :
    \begin{itemize}
        \item Le \textbf{client} est le client du restaurant qui consulte le menu et passe commande.
        \item Le \textbf{backend} est le chef de cuisine qui reçoit la commande, prépare le plat selon les recettes.
        \item La \textbf{base de données} est le garde-manger qui stocke les ingrédients et les recettes.
    \end{itemize}
\end{analogybox}

\section{Qu'est-ce qu'une API ?}

\begin{definitionbox}
    \textbf{API (Application Programming Interface)} : Ensemble de règles et de protocoles qui permettent à différentes applications logicielles de communiquer entre elles. Dans le contexte web, une API REST définit comment le client peut demander des données au serveur et comment le serveur doit répondre.
\end{definitionbox}

Une API REST (Representational State Transfer) utilise les méthodes HTTP standard :
\begin{itemize}
    \item \textbf{GET} : Récupérer des données (lecture)
    \item \textbf{POST} : Créer de nouvelles données (création)
    \item \textbf{PUT/PATCH} : Modifier des données existantes (mise à jour)
    \item \textbf{DELETE} : Supprimer des données (suppression)
\end{itemize}

\section{Pourquoi un Blog est un Projet Pédagogique Idéal ?}

Le projet de blog est particulièrement adapté à l'apprentissage du développement backend car :

\begin{itemize}
    \item Il implique des \textbf{relations complexes} entre les données (utilisateurs, articles, commentaires, tags).
    \item Il nécessite une \textbf{authentification} et une \textbf{autorisation} (rôles utilisateur).
    \item Il illustre parfaitement le \textbf{CRUD} (Create, Read, Update, Delete).
    \item Il permet de comprendre la \textbf{sécurité} (injection SQL, hachage de mots de passe).
    \item Il est suffisamment simple pour être compris rapidement, mais assez complexe pour être pédagogique.
\end{itemize}

\section{Architecture Complète du Projet}

Voici l'architecture complète de notre application :

\begin{center}
\begin{tikzpicture}[node distance=2cm, auto, thick]
    \tikzstyle{block} = [rectangle, draw, fill=deepblue!20, text width=3cm, text centered, rounded corners, minimum height=1.5cm]
    \tikzstyle{database} = [cylinder, draw, fill=green!20, shape border rotate=90, aspect=0.25, minimum height=2cm, minimum width=2cm, text centered]
    \tikzstyle{arrow} = [->, >=stealth, thick]
    
    \node[block] (browser) {Navigateur\\(HTML/CSS/JS)};
    \node[block, right=3cm of browser] (express) {Serveur Express\\(Node.js)};
    \node[database, right=3cm of express] (postgres) {PostgreSQL};
    
    \draw[arrow] (browser) -- node[above] {HTTP/JSON} (express);
    \draw[arrow] (express) -- node[above] {Requêtes SQL} (postgres);
    \draw[arrow] (postgres) -- node[below] {Données} (express);
    \draw[arrow] (express) -- node[below] {Réponses JSON} (browser);
    
    \node[below=0.5cm of browser] {Frontend};
    \node[below=0.5cm of express] {Backend};
    \node[below=0.5cm of postgres] {Base de données};
\end{tikzpicture}
\end{center}

\begin{architecturebox}
    \textbf{Flux de Requête Complet} :
    \begin{enumerate}
        \item L'utilisateur navigue sur le site dans son navigateur.
        \item Le navigateur envoie une requête HTTP au serveur Express.
        \item Express traite la requête via les middlewares et les routes.
        \item Le serveur interroge PostgreSQL via le driver pg.
        \item PostgreSQL retourne les données au serveur.
        \item Le serveur formate la réponse en JSON et l'envoie au navigateur.
        \item Le navigateur affiche les données à l'utilisateur.
    \end{enumerate}
\end{architecturebox}

% ============================================================================
% CHAPITRE 2 : INITIALISATION DU PROJET NODE.JS
% ============================================================================
\chapter{Initialisation du Projet Node.js}

\section{Installation de Node.js}

Node.js est un environnement d'exécution JavaScript côté serveur. Il permet d'exécuter du JavaScript en dehors du navigateur, ce qui est essentiel pour le développement backend.

\begin{definitionbox}
    \textbf{Node.js} : Environnement d'exécution JavaScript open-source, multiplateforme, basé sur le moteur V8 de Chrome, permettant d'exécuter du JavaScript côté serveur.
\end{definitionbox}

\subsection{Vérification de l'Installation}

Avant de commencer, vérifiez que Node.js est installé sur votre machine :

\begin{lstlisting}[style=bash, caption=Vérification de l'installation de Node.js]
node --version
npm --version
\end{lstlisting}

Si Node.js n'est pas installé, téléchargez-le depuis \url{https://nodejs.org/} et installez la version LTS (Long Term Support) recommandée.

\section{Création du Dossier du Projet}

Créez un nouveau dossier pour votre projet :

\begin{lstlisting}[style=bash, caption=Création du dossier du projet]
mkdir blog-multi-role
cd blog-multi-role
\end{lstlisting}

\section{Initialisation avec npm}

npm (Node Package Manager) est le gestionnaire de paquets de Node.js. Il permet de gérer les dépendances du projet.

\begin{lstlisting}[style=bash, caption=Initialisation du projet avec npm]
npm init -y
\end{lstlisting}

L'option \texttt{-y} permet de générer un fichier \texttt{package.json} avec les valeurs par défaut sans poser de questions interactives.

\subsection{Structure du Fichier package.json}

Le fichier \texttt{package.json} est le "carte d'identité" de votre projet. Il contient des métadonnées essentielles :

\begin{lstlisting}[style=javascript, caption=Exemple de package.json]
{
  "name": "blog-multi-role",
  "version": "1.0.0",
  "description": "Blog multi-rôle avec Node.js, Express et PostgreSQL",
  "main": "server.js",
  "scripts": {
    "start": "node server.js",
    "dev": "nodemon server.js"
  },
  "keywords": ["blog", "nodejs", "express", "postgresql"],
  "author": "Licence 1 Informatique",
  "license": "ISC",
  "dependencies": {
    "express": "^4.18.2",
    "pg": "^8.11.0",
    "bcrypt": "^5.1.0",
    "jsonwebtoken": "^9.0.0",
    "dotenv": "^16.0.3"
  },
  "devDependencies": {
    "nodemon": "^2.0.22"
  }
}
\end{lstlisting}

\begin{definitionbox}
    \textbf{Explication des champs} :
    \begin{itemize}
        \item \textbf{name} : Nom unique du projet
        \item \textbf{version} : Version du projet (suivant le sémantique versioning)
        \item \textbf{main} : Point d'entrée principal de l'application
        \item \textbf{scripts} : Commandes personnalisées pour exécuter le projet
        \item \textbf{dependencies} : Bibliothèques nécessaires en production
        \item \textbf{devDependencies} : Bibliothèques nécessaires uniquement en développement
    \end{itemize}
\end{definitionbox}

\subsection{Versioning Sémantique (SemVer)}

Le versioning sémantique suit le format \texttt{MAJEUR.MINEUR.CORRECTIF} :

\begin{itemize}
    \item \textbf{MAJEUR} : Changements incompatibles avec l'API précédente (ex: 1.0.0 → 2.0.0)
    \item \textbf{MINEUR} : Nouvelles fonctionnalités compatibles (ex: 1.0.0 → 1.1.0)
    \item \textbf{CORRECTIF} : Corrections de bugs compatibles (ex: 1.0.0 → 1.0.1)
\end{itemize}

\section{Structure Professionnelle du Projet}

Une structure de projet bien organisée est essentielle pour la maintenabilité et la collaboration.

\begin{lstlisting}[style=bash, caption=Structure complète du projet]
blog-multi-role/
├── server.js              # Point d'entrée du serveur
├── db.js                  # Configuration de la base de données
├── .env                   # Variables d'environnement (NE PAS COMMIT)
├── .gitignore             # Fichiers à ignorer par Git
├── package.json           # Configuration du projet
├── package-lock.json      # Verrouillage des versions
├── middleware/            # Middlewares Express
│   ├── auth.js           # Vérification JWT
│   └── roles.js          # Vérification des rôles
├── routes/                # Routes de l'API
│   ├── auth.js           # Routes d'authentification
│   ├── blogs.js          # Routes des blogs
│   └── commentaires.js   # Routes des commentaires
├── public/                # Fichiers statiques
│   ├── index.html        # Page d'accueil
│   ├── login.html        # Page de connexion
│   ├── register.html     # Page d'inscription
│   ├── css/
│   │   └── style.css     # Styles CSS
│   └── js/
│       └── app.js        # JavaScript frontend
├── sql/                   # Scripts SQL
│   └── schema.sql        # Schéma de la base de données
└── README.md              # Documentation du projet
\end{lstlisting}

\begin{architecturebox}
    \textbf{Explication des Dossiers} :
    \begin{itemize}
        \item \textbf{middleware/} : Contient les fonctions intermédiaires qui traitent les requêtes avant les routes (authentification, validation, etc.).
        \item \textbf{routes/} : Contient les définitions des routes de l'API, organisées par fonctionnalité.
        \item \textbf{public/} : Contient les fichiers statiques servis aux clients (HTML, CSS, JavaScript, images).
        \item \textbf{sql/} : Contient les scripts SQL pour la création et la manipulation de la base de données.
    \end{itemize}
\end{architecturebox}

\begin{bestpracticebox}
    \textbf{Principe de Séparation des Préoccupations} :
    Chaque dossier a une responsabilité unique et bien définie. Cette séparation rend le code plus maintenable, testable et compréhensible.
\end{bestpracticebox}

\end{document}
